\documentclass[12pt,a4paper]{article}
\usepackage[utf8]{inputenc}
\usepackage[korean]{babel}
\usepackage{amsmath,amssymb,amsthm}
\usepackage{geometry}
\usepackage{enumitem}

\geometry{margin=2.5cm}

\newtheorem{theorem}{정리}
\newtheorem{lemma}{보조정리}
\theoremstyle{definition}
\newtheorem{example}{반례}

\title{\textbf{조건부 거듭제곱 합 부등식}}
\author{}
\date{2025년 11월 13일}

\begin{document}

\maketitle

\begin{theorem}
\textbf{조건:} $n$개의 비음수 실수 $x_1, x_2, \ldots, x_n \geq 0$이
\[
\sum_{i=1}^n x_i = \sum_{i=1}^n x_i^2
\]
을 만족한다고 하자.

\textbf{결론:} 그러면 다음 부등식이 성립한다:
\[
\sum_{i=1}^n x_i^4 \geq \sum_{i=1}^n x_i^3
\]

\textbf{등호 조건:} 등호가 성립하는 것은 모든 $x_i \in \{0, 1\}$일 때이며 오직 그때뿐이다.

\textit{주의:} 음수를 허용하면 이 부등식은 거짓이다 (아래 반례 참조).
\end{theorem}

\section*{증명 1 --- 한 줄 인수분해}

\begin{proof}
조건 $\sum x_i = \sum x_i^2$로부터
\[
\sum_{i=1}^n (x_i^2 - x_i) = 0
\]

한편, 임의의 $x$에 대해 다음 항등식이 성립한다:
\[
(x^4 - x^3) - (x^2 - x) = x(x-1)(x^2-1)
\]

이제 $x \geq 0$일 때 $x(x-1)(x^2-1)$의 부호를 조사하자:
\begin{itemize}[leftmargin=1.5cm]
    \item $x \in [0, 1]$일 때: $x \geq 0$, $(x-1) \leq 0$, $(x^2-1) \leq 0$

    따라서 $x(x-1)(x^2-1) = x \cdot [(-) \times (-)] \geq 0$

    \item $x \geq 1$일 때: $x \geq 0$, $(x-1) \geq 0$, $(x^2-1) \geq 0$

    따라서 $x(x-1)(x^2-1) \geq 0$
\end{itemize}

결론적으로 모든 $x \geq 0$에 대해
\[
(x^4 - x^3) - (x^2 - x) \geq 0
\]

이제 각 $x_i$에 대해 위 부등식을 적용하고 합을 취하면:
\begin{align*}
\sum_{i=1}^n (x_i^4 - x_i^3) &= \sum_{i=1}^n \left[(x_i^4 - x_i^3) - (x_i^2 - x_i)\right] + \sum_{i=1}^n (x_i^2 - x_i) \\
&= \sum_{i=1}^n \left[(x_i^4 - x_i^3) - (x_i^2 - x_i)\right] + 0 \\
&\geq 0
\end{align*}

따라서 $\sum x_i^4 \geq \sum x_i^3$이다.

\textbf{등호 조건:}
등호가 성립하려면 각 항에서 $x_i(x_i-1)(x_i^2-1) = 0$이어야 한다.
$x_i \geq 0$이므로, 이는 $x_i = 0$ 또는 $x_i = 1$을 의미한다.
\end{proof}

\section*{증명 2 --- 단조성 비교 (대조항 이용)}

\begin{proof}
먼저 다음 보조 부등식을 증명한다:
\[
x^3(1-x) \leq x(1-x) \quad \text{for all } x \geq 0
\]

\textbf{보조 부등식의 증명:}
\begin{itemize}[leftmargin=1.5cm]
    \item $0 \leq x \leq 1$일 때: $1-x \geq 0$이고 $x^3 \leq x$이므로

    $x^3(1-x) \leq x(1-x)$

    \item $x \geq 1$일 때: $1-x \leq 0$이고 $x^3 \geq x$이므로

    $x^3(1-x) \leq x(1-x)$ (둘 다 음수이고 절댓값이 더 큼)
\end{itemize}

이제 주 부등식을 증명한다. 각 $x_i$에 대해 보조 부등식을 적용하면:
\[
\sum_{i=1}^n (x_i^3 - x_i^4) = \sum_{i=1}^n x_i^3(1-x_i) \leq \sum_{i=1}^n x_i(1-x_i)
\]

우변을 전개하면:
\[
\sum_{i=1}^n x_i(1-x_i) = \sum_{i=1}^n x_i - \sum_{i=1}^n x_i^2 = 0
\]
(조건에 의해 $\sum x_i = \sum x_i^2$)

따라서:
\[
\sum_{i=1}^n (x_i^3 - x_i^4) \leq 0
\]

즉, $\sum x_i^4 \geq \sum x_i^3$이다.

\textbf{등호 조건:}
등호가 성립하려면 각 $x_i$에서 $x_i^3(1-x_i) = x_i(1-x_i)$이어야 한다.
이는 $x_i(1-x_i)(x_i^2-1) = 0$을 의미하고, $x_i \geq 0$이므로 $x_i \in \{0, 1\}$이다.
\end{proof}

\section*{반례: 음수를 허용하면 부등식이 거짓}

\begin{example}
$n = 2$이고 $x_1 = \frac{3}{5}$, $x_2 = -\frac{1}{5}$라고 하자.

\textbf{조건 확인:}
\[
\sum_{i=1}^2 x_i = \frac{3}{5} - \frac{1}{5} = \frac{2}{5}
\]
\[
\sum_{i=1}^2 x_i^2 = \left(\frac{3}{5}\right)^2 + \left(-\frac{1}{5}\right)^2 = \frac{9}{25} + \frac{1}{25} = \frac{10}{25} = \frac{2}{5}
\]
조건 $\sum x_i = \sum x_i^2$가 성립한다. \checkmark

\textbf{부등식 검증:}
\[
\sum_{i=1}^2 x_i^3 = \left(\frac{3}{5}\right)^3 + \left(-\frac{1}{5}\right)^3 = \frac{27}{125} - \frac{1}{125} = \frac{26}{125}
\]
\[
\sum_{i=1}^2 x_i^4 = \left(\frac{3}{5}\right)^4 + \left(-\frac{1}{5}\right)^4 = \frac{81}{625} + \frac{1}{625} = \frac{82}{625}
\]

비교하면:
\[
\frac{26}{125} = \frac{130}{625} > \frac{82}{625}
\]

즉, $\sum x_i^3 > \sum x_i^4$이므로 부등식이 \textbf{역전}된다!

\textit{결론:} 비음수 조건 ($x_i \geq 0$)이 없으면 부등식은 일반적으로 성립하지 않는다.
\end{example}

\section*{요약}

이 정리는 다음을 보여준다:

\begin{enumerate}
    \item \textbf{비음수 조건의 필수성:}
    조건 $\sum x_i = \sum x_i^2$만으로는 부족하며, $x_i \geq 0$이 반드시 필요하다.

    \item \textbf{두 가지 증명 방법:}
    인수분해를 통한 직접 증명과 대조항을 이용한 비교 증명이 모두 가능하다.

    \item \textbf{등호 조건의 명확성:}
    등호는 모든 $x_i$가 0 또는 1일 때이며 오직 그때뿐이다.

    \item \textbf{반례의 존재:}
    음수를 허용하면 조건을 만족하면서도 부등식이 역전되는 경우가 존재한다.
\end{enumerate}

\end{document}
